%%%%%%%%%%%%%%%%%%%%%%%%%%%%%%%%%%%%%%%%%
% FRI Data Science_report LaTeX Template
% Version 1.0 (28/1/2020)
% 
% Jure Demšar (jure.demsar@fri.uni-lj.si)
%
% Based on MicromouseSymp article template by:
% Mathias Legrand (legrand.mathias@gmail.com) 
% With extensive modifications by:
% Antonio Valente (antonio.luis.valente@gmail.com)
%
% License:
% CC BY-NC-SA 3.0 (http://creativecommons.org/licenses/by-nc-sa/3.0/)
%
%%%%%%%%%%%%%%%%%%%%%%%%%%%%%%%%%%%%%%%%%


%----------------------------------------------------------------------------------------
%	PACKAGES AND OTHER DOCUMENT CONFIGURATIONS
%----------------------------------------------------------------------------------------
\documentclass[fleqn,moreauthors,10pt]{ds_report}
\usepackage[english]{babel}

\graphicspath{{fig/}}




%----------------------------------------------------------------------------------------
%	ARTICLE INFORMATION
%----------------------------------------------------------------------------------------

% Header
\JournalInfo{FRI Natural language processing course 2026}

% Interim or final report
\Archive{Project report} 
%\Archive{Final report} 

% Article title
\PaperTitle{Translation assistant for medical terminology} 

% Authors (student competitors) and their info
\Authors{Romina Mihalič, Anakhita Khademipur, Maša Pipan, Luka Škoda}

% Advisors
\affiliation{\textit{Advisors: Aleš Žagar}}

% Keywords
\Keywords{translation, Greek, English, terminology}
\newcommand{\keywordname}{Keywords}


%----------------------------------------------------------------------------------------
%	ABSTRACT
%----------------------------------------------------------------------------------------

\Abstract{

	For this project we are going to create an AI tranlsation assistant, that will be able to translate between Greek and English. We will use the T5 transformer architecture, which is a state-of-the-art model for natural language processing tasks. We will train the model on a large dataset of parallel sentences in Greek and English, and we will evaluate its performance on a test set of sentences that it has not seen before. We will also compare our model's performance to other translation models, such as Google Translate and DeepL. The model focuses on medical terminology.

}

%----------------------------------------------------------------------------------------

\begin{document}

% Makes all text pages the same height
\flushbottom 

% Print the title and abstract box
\maketitle 

% Removes page numbering from the first page
\thispagestyle{empty} 

%----------------------------------------------------------------------------------------
%	ARTICLE CONTENTS
%----------------------------------------------------------------------------------------

\section*{Introduction}

Accurate translation in specialized fields like medicine is difficult, especially when using general machine translation tools. While systems like Google Translate or large language models can produce fluent text, they often make mistakes with medical terminology or use inconsistent translations. This is a problem because in medical contexts, correct terminology is very important. 

In recent years, new approaches such as Retrieval-Augmented Generation (RAG) have been proposed to improve this. RAG combines information retrieval with text generation: instead of translating only based on the model, it first retrieves relevant documents and then uses them to produce a more accurate and grounded output.  

There are also many existing tools and frameworks for language processing, such as NLTK or spaCy, but they are general-purpose and do not focus on domain-specific translation.  

For this project, we focus on the medical domain and translation from English into Greek and Slovenian. We use resources from the CLARIN infrastructure, which provides access to many language corpora and tools for research.  

The main corpus used is the EMEA corpus, which contains medical texts from the European Medicines Agency and is suitable for extracting domain-specific terminology and context. In addition, we use another CLARIN corpus (ECDC-TM) to increase the amount of available data. For additional terminology grounding, we use a medical glossary, which includes English–Greek medical terms and helps ensure consistent translation. Such terminological resources are important because they standardize medical vocabulary and improve translation quality.  

The goal of this project is to build a translation assistant that produces more accurate and consistent medical translations by combining retrieved medical texts with glossary-based terminology.
%------------------------------------------------

\section*{Methods}

First of our proposed medical corpora is the EMEA Corpus which is a multilingual parallel corpus made out of PDF documents from the European Medicines Agency\cite{clarin_emea_2015}. Our second corpus is the ECDC Translation Memory which is a bilingual parallel corpus as it offers a collection of texts and their manual translations for multiple language pairs. It was provided by the European Centre for Disease Prevention and Control and it focuses on the domain of public health\cite{clarin_ecdc_tm_2015}. The third resource we intend to use is a bilingual Glossary of Medical Terms which was originally posted on the foreign languages website of the School of Medicine. It was created to help out students who had taken terminology courses\cite{clarin_medical_glossary_2016}. 


%----------------------------------------------------------------------------------------
%	REFERENCE LIST
%----------------------------------------------------------------------------------------
\bibliographystyle{unsrt}
\bibliography{report}
\end{document}
